\documentclass[11pt]{article}
\usepackage[margin=1in]{geometry}
\usepackage{amsmath,amssymb,booktabs}
\usepackage[hidelinks]{hyperref}

\title{CV Quantum Computing for a DV Practitioner:\\
A Compact Introduction to CV Formalism, CV Algorithms, and CV-DV Differential-Equation Solvers}
\author{}
\date{\today}

\begin{document}
\maketitle

\section*{Who this note is for}
This note is for readers comfortable with discrete-variable (DV, qubit) quantum computing who want a practical and mathematically grounded entry into continuous-variable (CV) methods and hybrid CV-DV algorithms, especially for differential equations (DE/PDE).

\section{From DV to CV in one page}
A DV register uses finite-dimensional Hilbert spaces (qubits). A CV mode uses an infinite-dimensional harmonic-oscillator Hilbert space.

\subsection*{Core operators}
For one CV mode, define quadratures $\hat q,\hat p$ with
\begin{equation}
[\hat q,\hat p]=i\hbar,
\end{equation}
and ladder operators
\begin{equation}
\hat a = \frac{1}{\sqrt{2\hbar}}(\hat q+i\hat p),
\qquad
\hat a^\dagger = \frac{1}{\sqrt{2\hbar}}(\hat q-i\hat p).
\end{equation}
In this repository, the simulator convention uses $\hbar=2$.

\subsection*{State families}
\begin{itemize}
\item Fock states $\{|n\rangle\}_{n\ge 0}$: oscillator number basis.
\item Coherent states $|\alpha\rangle$: displaced vacuum, closest analog to classical field amplitudes.
\item Squeezed states: reduced variance in one quadrature, increased in the conjugate quadrature.
\item Gaussian states: completely specified by first and second moments of quadratures.
\item Non-Gaussian states: needed for universal CV computation and many advanced algorithms.
\end{itemize}

\subsection*{DV$\leftrightarrow$CV mental map}
\begin{center}
\begin{tabular}{ll}
\toprule
DV notion & CV analog \\
\midrule
Qubit state vector & Wavefunction / Wigner function of a mode \\
Pauli generators & Quadrature generators and polynomial Hamiltonians \\
Single-qubit rotations & Phase-space displacements/rotations/squeezing \\
Measurement in $Z/X$ & Homodyne (quadrature) / heterodyne / photon counting \\
Finite basis cutoff & Fock truncation ($N_{\text{cut}}$) \\
\bottomrule
\end{tabular}
\end{center}

\section{Minimal CV formalism you should know}

\subsection{Gaussian unitaries (easy to implement, not universal alone)}
Typical single/multi-mode Gaussian gates:
\begin{align}
D(\alpha) &= e^{\alpha \hat a^\dagger-\alpha^*\hat a} \quad \text{(displacement)},\\
S(\zeta) &= e^{\frac{1}{2}(\zeta^*\hat a^2-\zeta \hat a^{\dagger 2})} \quad \text{(squeezing)},\\
R(\theta) &= e^{i\theta \hat a^\dagger\hat a} \quad \text{(phase rotation)},
\end{align}
and beamsplitters between modes.

Gaussian states under Gaussian gates and Gaussian measurements are classically tractable in many regimes; this is similar in spirit to Clifford stabilizer tractability in DV.

\subsection{What makes CV universal}
Gaussian operations alone are not universal. Add at least one non-Gaussian resource (e.g., cubic phase gate/state, photon counting, or encoded non-Gaussian ancilla such as GKP states), and universality is recovered.

\subsection{Noise and error correction in CV}
CV hardware naturally suffers shift-like errors in phase space (small random drifts in $q,p$). GKP encoding maps a logical qubit into one oscillator and converts small shift noise into correctable discrete syndrome information.

\section{CV algorithm landscape (for DV readers)}

\subsection{Category A: Gaussian-information protocols}
Teleportation, entanglement distribution, Gaussian channels, continuous-variable QKD, and Gaussian sensing/estimation. Foundational formalism is covered by major CV reviews.

\subsection{Category B: Measurement-based CV computing}
Large CV cluster states (often optical) plus adaptive measurements implement computations. Fault-tolerance thresholds are discussed in finite-squeezing settings (with concatenated DV-style coding ideas).

\subsection{Category C: Hybrid CV-DV algorithms}
A common pattern is: use CV modes as powerful analog/spectral ancillas while keeping part of the logic in qubits. This is especially natural for block-encoding-like constructions, nonunitary dynamics, and differential-equation solvers.

\section{CV-DV DE/PDE solvers: what the algorithm is doing}

\subsection{Problem class}
A standard target is linear dynamics
\begin{equation}
\partial_t u(t) = -A u(t),
\qquad
u(t)=e^{-tA}u(0),
\end{equation}
with $A$ often sparse/structured from PDE discretization.

\subsection{LCHS viewpoint}
Linear Combination of Hamiltonian Simulation (LCHS) rewrites nonunitary evolution into weighted combinations of unitary-like evolutions that can be queried/compiled on quantum hardware. In hybrid CV-DV implementations, a CV ancilla can encode the continuous kernel/integral structure with fewer explicit discretization steps than purely DV constructions.

\subsection{Why CV ancillas help}
A qumode naturally represents continuous spectral filters. This can reduce overhead for integral transforms or spectral weighting and improve compilation structure for certain nonunitary problems.

\subsection{Current practical compromise (important)}
In near-term implementations (including this project), one often uses a truncated Fock basis and, sometimes, an incoherent Fock-mixture approximation:
\begin{equation}
\rho_{\text{CV}} \approx \sum_n |C_n|^2 |n\rangle\langle n|,
\end{equation}
instead of a fully coherent superposition. This is backend-friendly but can produce mixed post-selected DV states and introduces model error distinct from Trotter/truncation effects.

\section{How this maps to your repository}

\subsection{Heat-equation target in this codebase}
The DV target is the finite-difference heat generator $T$ (2-qubit, 4 interior points), with solution
\begin{equation}
u_{\text{PDE}}(t)=e^{-\alpha tT}u(0).
\end{equation}
The CV side prepares LCHS-like coefficients $C_n(r_{\text{target}},r',\beta_k)$, then aggregates post-selected DV observables over Fock components.

\subsection{Parameters that matter most in practice}
\begin{itemize}
\item $r_{\text{target}}$: post-selection squeezing scale.
\item $r'$: preparation-basis squeezing (must satisfy $r'<r_{\text{target}}$).
\item $\beta_k$ (kernel hyperparameter): shapes spectral weighting in coefficient construction.
\item Fock cutoff / $n_{\text{dim}}$: truncation-runtime tradeoff.
\item Trotter steps: usually a smaller error source than state-prep/model mismatch in your runs.
\end{itemize}

\subsection{Metric priority (aligned with your project goals)}
For PDE solving, rank by:
\begin{enumerate}
\item relative PDE-vector error (primary),
\item post-selection probability (secondary feasibility),
\item fidelity only as a diagnostic (especially when output is mixed).
\end{enumerate}

\section{A practical learning path for DV users}
\begin{enumerate}
\item Learn CV phase-space formalism and Gaussian mechanics first (operators, covariance matrices, homodyne).
\item Learn what non-Gaussian resources buy you (universality and error correction).
\item Read LCHS/nonunitary simulation papers to understand DE solver scaling.
\item Revisit your code with this decomposition: \textit{algorithmic target error} + \textit{numerical discretization error} + \textit{CV-state-prep/model error} + \textit{hardware/backend error}.
\end{enumerate}

\section{References and reading map}
\subsection*{Foundational CV references}
\begin{enumerate}
\item S. Lloyd and S. L. Braunstein, \emph{Quantum Computation over Continuous Variables}, Phys. Rev. Lett. 82, 1784 (1999). DOI: \url{https://doi.org/10.1103/PhysRevLett.82.1784}
\item S. L. Braunstein and P. van Loock, \emph{Quantum information with continuous variables}, Rev. Mod. Phys. 77, 513 (2005). DOI: \url{https://doi.org/10.1103/RevModPhys.77.513}
\item C. Weedbrook \emph{et al.}, \emph{Gaussian quantum information}, Rev. Mod. Phys. 84, 621 (2012). DOI: \url{https://doi.org/10.1103/RevModPhys.84.621}
\item D. Gottesman, A. Kitaev, and J. Preskill, \emph{Encoding a qubit in an oscillator}, Phys. Rev. A 64, 012310 (2001). DOI: \url{https://doi.org/10.1103/PhysRevA.64.012310}
\item N. C. Menicucci, \emph{Fault-Tolerant Measurement-Based Quantum Computing with Continuous-Variable Cluster States}, Phys. Rev. Lett. 112, 120504 (2014). DOI: \url{https://doi.org/10.1103/PhysRevLett.112.120504}
\end{enumerate}

\subsection*{DE/PDE and CV-DV algorithm references (including your res\_refs)}
\begin{enumerate}
\item D. An, J.-P. Liu, and L. Lin, \emph{Linear Combination of Hamiltonian Simulation for Nonunitary Dynamics with Optimal State Preparation Cost}, Phys. Rev. Lett. 131, 150603 (2023). DOI: \url{https://doi.org/10.1103/PhysRevLett.131.150603}
\item D. An, A. M. Childs, and L. Lin, \emph{Quantum algorithm for linear non-unitary dynamics with near-optimal dependence on all parameters} (2024). (file in \texttt{res\_refs})
\item M. Pocrnic, P. D. Johnson, A. Katabarwa, and N. Wiebe, \emph{Constant-Factor Improvements in Quantum Algorithms for Linear Differential Equations} (2024). (file in \texttt{res\_refs})
\item S. Jin and N. Liu, \emph{Analog quantum simulation of partial differential equations}, Quantum Sci. Technol. 9, 035047 (2024). DOI: \url{https://doi.org/10.1088/2058-9565/ad49cf}
\item L. Bell \emph{et al.}, \emph{Co-Designing Eigen- and Singular-value Transformation Oracles: From Algorithmic Applications to Hardware Compilation} (2025). (file \texttt{res\_refs/Contiuous LCU.pdf})
\item Y. Gan \emph{et al.}, \emph{Provably Efficient Quantum Algorithms for Solving Nonlinear Differential Equations Using Multiple Bosonic Modes Coupled with Qubits} (2025). arXiv:2511.09939.
\end{enumerate}

\subsection*{Repository-local context}
\begin{itemize}
\item See \texttt{heat\_eq\_postselect.py} for hybrid circuit structure and CV/DV operator mapping used in this project.
\item See \texttt{heat\_eq\_sensitivity\_refine.py} and \texttt{heat\_eq\_systematic\_optimize.py} for PDE-priority parameter search.
\end{itemize}

\end{document}
